% Paper 1 — deterministic (Phases 1–5). DRAFT v0.1 (2026-07-08).
% Build: cd paper/draft && tectonic main.tex
% NOTE: written in plain article class for drafting; port to the target
% journal template (JOSA A / Opt. Express / PRA) at submission.
% Citations marked [V] are verified against the source (paper/literature/
% LITERATURE_SWEEP.md); those marked TODO-verify must be checked before
% submission (see docs/LITERATURE.md caveats).
\documentclass[11pt]{article}
\usepackage[margin=2.4cm]{geometry}
\usepackage{graphicx,amsmath,amssymb,booktabs,xcolor}
\usepackage[hidelinks]{hyperref}
\graphicspath{{../../output/05_pulsed_pupils/}{../../output/03_annular/}%
{../../output/04_axicon/}{../../output/02_na_apodization/}{../../output/00_overview/}}
\newcommand{\todo}[1]{\textcolor{red}{[TODO: #1]}}
\setlength{\parskip}{2pt}

\title{Few-cycle pulses focused through annular and conical pupils at high
numerical aperture: bandwidth-invariant depth-of-focus extension and the
erosion of transverse structure}
\author{R.~A.~Romero\\ \small National Institute of Physics, University of the
Philippines Diliman\\ \small \todo{author list / affiliations to confirm}}
\date{Draft v0.1 --- \today}

\begin{document}
\maketitle

\begin{abstract}
Structured pupils --- annular apertures and conical (axicon) phases --- extend
the depth of focus (DOF) of a tightly focused beam, while few-cycle pulses
carry an octave-scale bandwidth that smears the fine structure of the focal
field. What happens when the two are combined at high numerical aperture (NA),
where the field is fully vectorial, has not been reported. Using
Richards--Wolf vector diffraction with full-spectrum temporal integration, we
show that the fractional DOF extension produced by an annular pupil,
$\mathrm{DOF}(\varepsilon)/\mathrm{DOF}(0)$, is \emph{invariant to few-cycle
bandwidth} --- constant to three digits from continuous wave down to a 1-fs
pulse ($44.35\!\to\!44.33$ at obstruction $\varepsilon=0.99$, $\mathrm{NA}=0.8$,
$\lambda_c=750$\,nm) and across apertures $\mathrm{NA}=0.6$--$1.0$ --- while the
transverse Bessel ring structure erodes \emph{non-monotonically} in pulse
width: the Linfoot structural content first dips below unity (moderate pulses
wash out the rings) before rising above it (few-cycle broadening of the central
lobe). A core/ring decomposition explains the non-monotonicity quantitatively.
Supporting results sharpen the deterministic toolkit: scalar-validity
thresholds at high NA are shown to be strongly anisotropic (dependent on
polarization, observation cut, and direction --- a single ``NA threshold'' is
ill-posed); Gaussian apodization \emph{extends} scalar validity; the
thin-annulus Bessel core is vector-deformed at high NA while the axicon's
apparent deviation from the Durnin $J_0^2$ profile is shown to be
stationary-phase error, vanishing as $1/\beta$. All results are validated
against analytic limits, published anchors, and first-principles
Maxwell/Helmholtz residual checks. The bandwidth-invariant DOF extension
identifies pulsed annular/Bessel illumination as a robust route to extended-
depth nonlinear microscopy with few-cycle sources.
\end{abstract}

\section{Introduction}

Tight focusing of ultrashort pulses underlies multiphoton microscopy,
micromachining, and strong-field light--matter interaction. At high numerical
aperture the focal field must be treated vectorially; the Richards--Wolf (RW)
diffraction integrals~\cite{Wolf1959,RichardsWolf1959} remain the standard
description, and their extension to broadband pulses --- summing monochromatic
vector fields incoherently over the pulse power spectrum --- was analyzed for
unstructured (clear circular) pupils by Romallosa, Bantang and
Saloma~\cite{Romallosa2003} [V], who quantified the smearing of the focal
structure of a 750-nm beam at pulse widths down to 1\,fs via Linfoot's image
criteria [V].

Independently, \emph{pupil structuring} is a mature route to focal
engineering. Annular apertures narrow the central lobe and extend the depth of
focus~\cite{SheppardWilson1978,Sheppard1988}; in the thin-ring limit the focal
field approaches Durnin's nondiffracting Bessel beam~\cite{Durnin1987}; conical
(axicon) phases produce an extended focal segment with a position-dependent
Bessel profile~\cite{McLeod1954,Durnin1987}. At high NA these structured foci
acquire vector character: double-ring and cylindrical-vector illumination
produce sub-wavelength needles of controlled polarization over many
wavelengths~\cite{Wang2010needle,Guo2014needle} [V]. On the application side,
axicon-generated Bessel--Gauss beams extend the two-photon depth of field by an
order of magnitude or more~\cite{Betzig2014Bessel} [V-mechanism], and Bessel
light sheets image deep into scattering tissue~\cite{Zhao2014BesselSIM}
\todo{verify exact numbers before submission}.

What is missing between these two literatures is their \emph{combination}: how
do structured-pupil foci respond to few-cycle bandwidth at high NA? The
bandwidth couples to the pupil geometry nontrivially --- every spectral
component sees the same geometric pupil but produces a focal pattern whose
scale is proportional to its wavelength, so spectral summation washes out
wavelength-scale structure (the Romallosa effect~\cite{Romallosa2003}) while
potentially preserving geometric ratios. To our knowledge no published work
reports the pulsed behavior of annular or axicon foci in the vectorial,
few-cycle regime (Sec.~\ref{sec:lit} discusses the closest prior art).

Here we answer that question with a validated deterministic toolkit
(Sec.~\ref{sec:methods}) and report four results (Sec.~\ref{sec:results}):
(i)~scalar-validity thresholds at high NA are strongly \emph{anisotropic} ---
they depend on input polarization, on the observation cut, and on direction
(transverse vs axial), so a single scalar-breakdown NA is ill-posed; Gaussian
apodization systematically \emph{extends} scalar validity;
(ii)~the annular DOF law $\mathrm{DOF}(\varepsilon)/\mathrm{DOF}(0)
=1/(1-\varepsilon^2)$ holds to three digits paraxially but falls ${\sim}40\%$
short at $\sin\alpha=0.9$, and the thin-annulus Bessel core is genuinely
vector-deformed at high NA (deviation 0.163 from $J_0^2$ vs 0.004 at low NA);
(iii)~the axicon's transverse profile deviation from the predicted local
$J_0^2$ is a stationary-phase artifact that decays as $1/\beta$ (0.134 at
$\beta{=}5$ to 0.006 at $\beta{=}2560$), not a vector effect --- with practical
validity ceilings; and
(iv)~\emph{the headline}: the fractional annular DOF extension is invariant to
few-cycle bandwidth across $\mathrm{NA}=0.6$--$1.0$, while the transverse ring
contrast erodes non-monotonically in pulse width, an effect we explain
quantitatively by decomposing the pattern into its Bessel core and rings.

\section{Theory and methods}\label{sec:methods}

\subsection{Vector focal fields}
For an aplanatic system of semi-aperture $\alpha$ ($\sin\alpha=\mathrm{NA}/n$)
illuminated by an $x$-polarized beam with pupil amplitude $A(\theta)$, the
focal field is given by the RW integrals~\cite{RichardsWolf1959}
\begin{equation}
I_m(r,z;k)=\int_0^{\alpha}\! A(\theta)\sqrt{\cos\theta}\,
K_m(\theta)\, J_m(kr\sin\theta)\, e^{\mathrm{i}kz\cos\theta}\,\mathrm{d}\theta,
\qquad m=0,1,2,
\label{eq:rw}
\end{equation}
with kernels $K_0=\sin\theta(1{+}\cos\theta)$, $K_1=\sin^2\theta$,
$K_2=\sin\theta(1{-}\cos\theta)$, and field components
$\mathbf{E}\propto\big(I_0+I_2\cos2\phi,\; I_2\sin2\phi,\;
-2\mathrm{i}I_1\cos\phi\big)$. Following Ref.~\cite{Romallosa2003} we quote the
meridional cut with true zeros, $|I_0-I_2|^2$ (the $y$--$z$ plane), and verify
key results on the orthogonal cut $|I_0{+}I_2|^2{+}4|I_1|^2$ and the azimuthal
average.

Annular pupils of obstruction $\varepsilon$ are treated by field-level Babinet
subtraction, $I_m^{\mathrm{ann}}=I_m(\alpha)-I_m(\alpha_{\mathrm{in}})$ with
$\sin\alpha_{\mathrm{in}}=\varepsilon\sin\alpha$; apodization uses the
truncated-Gaussian $A(\theta)=\exp(-\alpha_t\sin^2\theta/\sin^2\alpha)$ of
Horv\'ath and Bor~\cite{HorvathBor2003}; the axicon is the conical pupil phase
$A(\theta)=\exp(\mathrm{i}\beta\sin\theta/\sin\alpha)$.

\subsection{Few-cycle pulses}
An unchirped Gaussian pulse of intensity FWHM $\tau$ has power spectrum
$|E(\omega)|^2\propto\exp[-(\omega-\omega_c)^2/2a]$ with $a=2\ln2/\tau^2$
(spectral FWHM $4\ln2/\tau$; the transform-limited product
$\Delta\nu\,\Delta t=2\ln2/\pi$). The time-integrated focal intensity is the
incoherent spectral sum~\cite{Romallosa2003}
\begin{equation}
I(\mathbf{r})=\int_0^{2\omega_c}\! |E(\omega)|^2\,
\big|\mathbf{E}(\mathbf{r};\omega)\big|^2\,\mathrm{d}\omega .
\label{eq:eq10}
\end{equation}
The integration window matters: we find (Appendix~\ref{app:window}) that
Ref.~\cite{Romallosa2003}'s published anchors are reproduced only when
Eq.~\eqref{eq:eq10} is evaluated over the \emph{full positive spectral support}
$\omega\in(0,2\omega_c)$; restricting to the FWHM band quoted in the text
discards 24\% of the spectral weight and fails every anchor.

\subsection{Metrics and validation}
Pulsed-vs-cw pattern change is quantified by Linfoot's criteria (fidelity $F$,
structural content $S$, correlation quality $Q$; $2Q-S=F$) on peak-normalized
profiles~\cite{Romallosa2003,Linfoot}. DOF denotes the axial FWHM of the
on-axis intensity in optical units $u=k\sin^2\!\alpha\,z$.

Every layer of the computation is validated before reuse: low-NA limits
reproduce the Airy and Durnin $J_0^2$ patterns to $5\times10^{-3}$; Babinet at
$\varepsilon{=}0$ is exact to $10^{-12}$; the published anchors of
Ref.~\cite{Romallosa2003} (axial $S=1.058$ vs.\ 1.057; transverse FWHM
472.8\,nm; zero rings at 0.555 and 3.05\,$\mu$m) and the annular closed forms
of Ref.~\cite{SheppardWilson1978} (0.8\%) are reproduced. Beyond benchmark
matching, the implementation passes \emph{first-principles} residual checks:
the scalar field satisfies the Helmholtz equation and the vector field
satisfies $\nabla\!\cdot\!\mathbf{E}=0$ at the finite-difference floor
($\sim\!10^{-4}$), with deliberately broken controls (wrong $k$; $E_z$ kernel
factor) failing at the $10^{-1}$ level --- pinning the kernels of
Eq.~\eqref{eq:rw} to Maxwell's equations independently of any literature
anchor. A 40-test regression suite locks these results.
\todo{cite code/data repository (Zenodo DOI) at submission}

\section{Results}\label{sec:results}

\subsection{Anisotropic scalar-validity thresholds; apodization extends
scalar validity}\label{sec:thresholds}
Comparing vector and scalar predictions of the focal FWHM as NA grows shows
that the scalar-breakdown threshold is not a single number
(Table~\ref{tab:thresh}): at the 1\% level the worst case (transverse,
linear-$x$ input, cut parallel to the polarization) fails at
$\mathrm{NA}\approx0.30$, while the perpendicular cut holds to ${\approx}0.85$
and the axial profile to ${\approx}1.15$. The anisotropy \emph{is} the
longitudinal field: $E_z$ fills the parallel-cut zeros and elongates the spot
along the polarization axis. Gaussian underfilling ($\alpha_t{=}4$) de-weights
marginal rays and \emph{extends} the worst-case 1\%/5\% thresholds from
0.30/0.60 to 0.40/0.85.

\begin{table}[htb]\centering\small
\caption{NA at which the scalar prediction of the focal FWHM deviates by
1\%/5\% (uniform pupil, $\lambda_c=750$\,nm, $n=1.3$).}
\label{tab:thresh}
\begin{tabular}{lcc}\toprule
quantity & 1\% & 5\% \\ \midrule
transverse, linear-$x$, $\parallel$ cut (worst) & 0.30 & 0.60\\
transverse, circular input & 0.45 & 0.90\\
transverse, linear-$x$, $\perp$ cut & 0.85 & ---\\
axial (any polarization) & 1.15 & ---\\ \bottomrule
\end{tabular}
\end{table}

\subsection{Annular apertures: DOF law and the vector Bessel core}
The paraxial law $\mathrm{DOF}(\varepsilon)/\mathrm{DOF}(0)=1/(1-\varepsilon^2)$
holds to three digits at low NA (50.13 vs.\ 50.25 at $\varepsilon=0.99$) and
falls ${\sim}40\%$ short at $\sin\alpha=0.9$ (30.4). The Strehl cost follows
$(1-\varepsilon^2)^2$. In the thin-ring limit the transverse profile approaches
$J_0^2(v)$ with deviation 0.004 at low NA but \textbf{0.163} at
$\sin\alpha=0.9$ --- a clean, isolated \emph{vector} deformation of the Bessel
core (the single-cone geometry isolates the $(1\pm\cos\theta)$ kernels at one
angle).

\subsection{The axicon deviation is stationary phase, not vector}
The axicon focal segment carries a local Bessel profile
$J_0^2(v\sin\theta^*/\sin\alpha)$ at the stationary cone angle
$\theta^*(u)$. Its 4--8\% deviation from this prediction is \emph{equal} at low
and high NA, and a $\beta$-scan shows it decays as $1/\beta$ (0.134 at
$\beta{=}5$ to 0.006 at $\beta{=}2560$): it is the stationary-phase
approximation error of the analytic model, not a failure of the field
computation (which matches the full Debye integral to $10^{-3}$ throughout).
Practical validity ceilings: the pupil-phase sampling resolves
$\beta\gtrsim2500$, but for a physical objective ($f\approx2$\,mm) the focal
segment leaves the Debye focal region at $\beta\approx160$.

\subsection{Headline: pulsed annular pupils --- invariant DOF, eroding
rings}\label{sec:headline}
Figure~\ref{fig:headline} combines the validated spectral sum
[Eq.~\eqref{eq:eq10}] with the validated annular fields. Two phenomena
coexist:

\textbf{(a) The DOF extension is bandwidth-invariant.} At $\mathrm{NA}=0.8$,
$\mathrm{DOF}(\varepsilon)/\mathrm{DOF}(0)$ changes from 44.35 (cw) to 44.33
(1\,fs) at $\varepsilon=0.99$ --- invariance to three digits at a bandwidth
where the transverse pattern is strongly smeared. The robustness map
(Fig.~\ref{fig:robust}) shows the same invariance across apertures: over
$\varepsilon=0.5$--$0.99$ and $\mathrm{NA}=0.6$--$1.0$ the cw$\to$1-fs spread
of the ratio never exceeds \textbf{0.10\%} (e.g.\ $47.30\!\to\!47.30$ at
$\mathrm{NA}{=}0.6$, $44.35\!\to\!44.33$ at $0.8$, $39.16\!\to\!39.12$ at
$1.0$, all at $\varepsilon=0.99$), with the largest spreads at the highest
aperture --- consistent with the weak residual NA-dependence of the ratio
itself.
The mechanism is geometric: each spectral component produces the same annular
cone geometry, so its \emph{fractional} DOF extension is
wavelength-independent; the spectral sum then averages nearly identical
ratios. The absolute DOF of each configuration lengthens only ${\sim}1\%$ from
cw to 1\,fs.

\textbf{(b) The transverse ring structure erodes non-monotonically.} The
Linfoot structural content of the transverse pattern (pulsed vs.\ cw at the
same $\varepsilon$) dips \emph{below} unity at moderate pulse widths
($S=0.988/0.987$ at 5/2\,fs, $\varepsilon=0.9$) before rising above it at
1\,fs ($S=1.058$). A core/ring decomposition at the first cw zero explains
this quantitatively: the Bessel-like central lobe is robust
($S_{\mathrm{core}}$ rises slowly, $1.00\!\to\!1.04$) while the rings first
wash out ($S_{\mathrm{ring}}$ dips to 0.72 --- rings of different spectral
components sit at different radii and average destructively) and are then
overfilled by the broadened central skirt at 1\,fs
($S_{\mathrm{ring}}=1.37$). The rings carry only 7\% of the pattern weight,
but with $S_{\mathrm{core}}\approx1$ the small ring deficit controls the sign
of $S-1$. The effect is robust: it is smooth on a dense $\tau$ grid, flat
against spectral sampling ($N_{\mathrm{freq}}=51$--$801$), and present in the
$\parallel$ cut, the $\perp$ cut, and the full azimuthally averaged vector
intensity.

\begin{figure}[htb]\centering
\includegraphics[width=0.95\textwidth]{verify_pulsed_annular.png}
\caption{Pulsed annular focusing at $\mathrm{NA}=0.8$: transverse ring filling
vs.\ $\tau$ at $\varepsilon=0.99$; full-disk reference; persistence of the
extended axial focus at 1\,fs; and the DOF ratio following the paraxial law
for all $\tau$. \todo{replace with publication-quality re-render}}
\label{fig:headline}
\end{figure}

\begin{figure}[htb]\centering
\includegraphics[width=0.95\textwidth]{paper_robustness.png}
\caption{Robustness of the headline: the normalized DOF ratio
$[\mathrm{DOF}(\varepsilon)/\mathrm{DOF}(0)](\tau)$ relative to its cw value,
for $\varepsilon=0.5,0.9,0.99$ at $\mathrm{NA}=0.6,0.8,1.0$.}
\label{fig:robust}
\end{figure}

\section{Relation to prior work}\label{sec:lit}
Ref.~\cite{Romallosa2003} treats few-cycle tight focusing for the clear
pupil; the annular/Bessel DOF literature
\cite{SheppardWilson1978,Durnin1987,Betzig2014Bessel} is monochromatic or
narrowband; high-NA vector needle work
\cite{Wang2010needle,Guo2014needle} is cw. Broadband extensions of Debye
focusing theory exist \cite{Botcherby2006matrix} [V] but do not address
structured pupils' pulsed response. We are not aware of a prior report of the
bandwidth-invariance of the annular DOF extension or of the non-monotonic ring
erosion. \todo{final targeted search (``pulsed annular depth of focus'',
``few-cycle Bessel focus'') before submission}

\section{Discussion and outlook}
For applications, the message is enabling: the axial benefit of
annular/Bessel illumination survives the broadest usable bandwidths, so
extended-DOF nonlinear microscopy \cite{Betzig2014Bessel} can employ few-cycle
sources without sacrificing the depth advantage --- what is traded is
transverse ring contrast, and in two-photon operation the rings are already
suppressed by the intensity-squared response. The anisotropic thresholds
(Sec.~\ref{sec:thresholds}) give practitioners a lookup for when scalar
modeling suffices, per polarization and per observable. The validated
deterministic fields constructed here are the source terms for a companion
study of these structured pulsed foci propagating through scattering media
(in preparation), where the wavelength dependence of scattering couples the
spectral and spatial degrees of freedom.

\appendix
\section{The spectral integration window of Ref.~[Romallosa 2003]}
\label{app:window}
Integrating Eq.~\eqref{eq:eq10} over the FWHM band
(483\,nm--1.67\,$\mu$m as quoted in the text of Ref.~\cite{Romallosa2003})
discards 24\% of the spectral weight at $\tau=1$\,fs and yields transverse
$S=1.020$, inconsistent with the published 1.062. Integrating the full
positive support $\omega\in(0,2\omega_c)$ reproduces every anchor (axial
$S=1.058$ vs.\ published 1.057; $F\approx0.96$; the pulsed pedestal
${\approx}0.05$ at $r=1\,\mu$m of their Fig.~2). We conclude the published
band describes the FWHM, not the integration window, and recommend the full
window for reproducibility.

\begin{thebibliography}{99}\small
\bibitem{Wolf1959} E. Wolf, Proc. R. Soc. Lond. A \textbf{253}, 349 (1959).
\bibitem{RichardsWolf1959} B. Richards and E. Wolf, Proc. R. Soc. Lond. A
\textbf{253}, 358 (1959).
\bibitem{Romallosa2003} K. M. Romallosa, J. Bantang, and C. Saloma, Phys. Rev.
A \textbf{68}, 033812 (2003). [V]
\bibitem{SheppardWilson1978} C. J. R. Sheppard and T. Wilson, Microwaves Opt.
Acoust. \textbf{2}, 105 (1978). \todo{verify}
\bibitem{Sheppard1988} C. J. R. Sheppard, annular-pupil DOF literature.
\todo{complete citation}
\bibitem{Durnin1987} J. Durnin, J. Opt. Soc. Am. A \textbf{4}, 651 (1987).
\bibitem{McLeod1954} J. H. McLeod, J. Opt. Soc. Am. \textbf{44}, 592 (1954).
\todo{verify}
\bibitem{Wang2010needle} double-ring radially polarized needle, Opt. Commun.
(2010), doi:S003039921000280X. [V]
\bibitem{Guo2014needle} double-ring azimuthally polarized needle, Opt. Lasers
Eng. \textbf{59}, 93 (2014). [V]
\bibitem{Betzig2014Bessel} axicon Bessel--Gauss two-photon microscopy
(Betzig group), PMC4032997. [V-mechanism] \todo{formal citation}
\bibitem{Zhao2014BesselSIM} two-photon Bessel light-sheet SIM, PMC4026892.
\todo{verify + formal citation}
\bibitem{HorvathBor2003} Z. L. Horv\'ath and Z. Bor, Opt. Commun. \textbf{222},
51 (2003). \todo{verify}
\bibitem{Botcherby2006matrix} matrix Debye integral for broadband focusing,
arXiv:physics/0609030 (2006). [V]
\bibitem{Linfoot} E. H. Linfoot, \emph{Fourier Methods in Optical Image
Evaluation} (Focal Press, 1964). \todo{verify}
\end{thebibliography}

\end{document}
